\subsection{What the two signs do and do not classify}

\begin{statement}{B.1}{Proved}{The precise finite-frame character classification}
Among continuous characters of $G=\GZ$ with values in $\Ftwo$
which vanish on $G_{\Q}$ and whose restriction to $\HH$ factors
through the two corrected finite frames $C_6\times C_4$, the
only characters are zero and $\Delta$.
\end{statement}
\begin{proof}
The full frame image and the Galois image are respectively
$C_6\times C_4$ and its index-two diagonal-parity subgroup,
as proved in B.23 below. A homomorphism from their quotient
to $\Ftwo$ is either zero or its unique nonzero character.
Thus the restriction is zero or $\delta$. The decomposition
$G=\HH G_{\Q}$ then uniquely determines a character
vanishing on $G_{\Q}$. This proves the stated restricted
classification, and makes no claim that every relative
character factors through these frames.
\end{proof}

\begin{statement}{B.2}{Proved}{The exact meaning of the constant sign function}
On $\HH$, let $\ell_h$ be the function on the two-element set
$\{j=0,j=1728\}$ with values
$((-1)^{a_3(h)},(-1)^{a_2(h)})$. Every such function occurs,
and it is constant exactly on $\ker\delta$.
On all of $G$, the gauge-invariant version is its class modulo
the common sign; its zero locus is exactly $\ker\Delta$.
This is the largest closed subgroup on which that obstruction
vanishes. No analogous statement about all five-point special
points follows without constructing their coordinates.
\end{statement}
\begin{proof}
R2 gives the pair $(1,0)$, while A.3 supplies a Galois pair
$(1,1)$. The image of the homomorphism $(a_3,a_2)$ therefore
contains all four elements of $\Ftwo^2$. Equality of the two
signs is the equation $a_3+a_2=0$. A subgroup on which the
equation holds is contained in its kernel, which is closed
and itself has the property. A.2--A.3 prove the asserted
all-multiplier formulation and explain why its two separate
components cannot be extended as abelian characters.
\end{proof}

\begin{statement}{B.3}{Proved}{A tower-action module carrying the sign line}
There is a finite-rank free $\Z_\ell$-module, for every odd
prime $\ell$, carrying a continuous $G$ action through its
action on a finite characteristic quotient of $\B$, and
containing a line with character $(-1)^\Delta$.
This construction is a permutation module of finite
automorphism residues; it is not a proof that a specified
level-12 homology module has such a line.
\end{statement}
\begin{proof}
Choose the finite characteristic level of A.10, let $Q$
be the finite image of $G$ in its outer automorphism group,
and let $\bar\Delta:Q\to\Ftwo$ be the induced character.
On $M=\Z_\ell[Q]$ use left translation of the basis $e_q$.
Then
\[
 v=\sum_{q\in Q}(-1)^{\bar\Delta(q)}e_q
\]
is nonzero, primitive, and satisfies
$g v=(-1)^{\Delta(g)}v$. The equality follows by replacing
$q$ with $g^{-1}q$ in the sum. The line is a direct
$\Z_\ell$-summand because one of its coordinates is a unit.
All maps factor through the specified finite tower action.
\end{proof}

\begin{statement}{B.4}{Refuted}{Not every module is distinguished from its sign twist}
Every $G$-module and its $\Delta$ twist restrict to canonically
identical $G_{\Q}$ actions. They need not be nonisomorphic
as $G$-modules, even when they contain the sign line.
\end{statement}
\begin{proof}
Let $\chi=(-1)^\Delta$ and take
$M=\Z_\ell\oplus\Z_\ell(\chi)$. The twisted module is
$\Z_\ell(\chi)\oplus\Z_\ell$, and interchanging the
summands is a $G$-equivariant isomorphism. On an element
with $\Delta=1$ the matrices are
\[
 D=\begin{pmatrix}1&0\\0&-1\end{pmatrix},\quad
 -D,\qquad
 P=\begin{pmatrix}0&1\\1&0\end{pmatrix},\quad PD=(-D)P.
\]
The executable algebra check verifies this identity.
On $G_{\Q}$ the twist is literally one, so restriction
does not change any action. In contrast the rank-one
trivial representation and its twist are nonisomorphic
over $\Z_\ell$ for odd $\ell$, since a putative
isomorphism would intertwine $1$ and $-1$ on the R2
element. This is the precise true replacement for the
universal nonisomorphism assertion.
\end{proof}

\begin{statement}{B.5}{Proved}{The powerset semiring statement, with its operations specified}
Let $\mathscr C(G)$ be the set of closed subsets of the
profinite group $G$, with addition given by union and
multiplication by setwise group product. Then
$S\mapsto\Delta(S)$ is a surjective homomorphism of
idempotent semirings to the powerset semiring of $\Ftwo$.
The image from closed subsets of $G_{\Q}$ is
$\{\varnothing,\{0\}\}$. The same holds for $\GT$
if and only if $\Delta|_{\GT}=0$.
\end{statement}
\begin{proof}
Finite unions of closed sets are closed, and products
are compact images of products of compact spaces, hence
closed. Group multiplication distributes over unions.
The zero is $\varnothing$ and the unit is $\{1\}$.
Direct image preserves unions, and the homomorphism law
gives $\Delta(ST)=\Delta(S)+\Delta(T)$ as sumsets.
The empty set, singleton identity, singleton R2 element,
and their two-element union realize all four subsets of
$\Ftwo$. Galois vanishing proves its asserted image.
The last equivalence follows by using singleton subsets.
If ``semiring'' were instead meant to use intersection
as multiplication, the assertion would be false: the
disjoint sets $\{1\}$ and $\{c\}$ both have image
$\{0\}$. Thus the operations matter.
\end{proof}

\begin{statement}{B.6}{Proved}{Continuous ring maps are precisely continuous unit characters}
For odd $\ell$, continuous unital ring homomorphisms
\[
 \Z_\ell[[G]]\longrightarrow\Z_\ell
\]
which agree with augmentation on $\Z_\ell[[G_{\Q}]]$
are in bijection with continuous characters
$\chi:G\to\Z_\ell^\times$ trivial on $G_{\Q}$.
If $K_G$ is the relative abelianization in P.4 below,
this set is
\[
 \Hom_{\rm cont}(K_G,\mu_{\ell-1})
 \times\Hom_{\rm cont}(K_G,\Z_\ell).
\]
In particular augmentation and $g\mapsto(-1)^{\Delta(g)}$
are distinct such maps. This is an exact classification
in terms of $K_G$, whose independent computation is
still required to list all characters explicitly.
\end{statement}
\begin{proof}
A unital continuous map fixes the ordinary integers
and therefore fixes $\Z_\ell$ by density. Group
elements are units and give a continuous character.
Conversely, reduction of a continuous unit character
modulo $\ell^n$ has finite image and open kernel.
Extending linearly on the corresponding finite group
ring and passing to the inverse limit gives the ring
map; density of the group algebra proves uniqueness.
Triviality on $G_{\Q}$ is precisely the augmentation
condition. Factorization through the abelianization
and its split cyclotomic factor gives $K_G$.
Finally
$\Z_\ell^\times=\mu_{\ell-1}\times(1+\ell\Z_\ell)$
and the convergent $\ell$-adic logarithm identifies
the second factor with the additive group $\ell\Z_\ell$,
which is isomorphic to $\Z_\ell$. The R2 value makes
the two displayed maps distinct.
\end{proof}

\begin{statement}{B.7}{Proved}{The relative residual tower detects every element}
For cofinal characteristic open subgroups $U$ of $\B$,
\[
 \B\simeq\varprojlim_U\B/R_{\rm nil}(U).
\]
An automorphism compatible with every quotient and
every transition map is therefore an actual automorphism
of $\B$. Reaching each quotient separately does not
produce a compatible system of automorphisms satisfying
an additional equation such as the pentagon.
\end{statement}
\begin{proof}
If $V\subset U$, restriction of a finite nilpotent
quotient of $U$ shows
$R_{\rm nil}(V)\subset R_{\rm nil}(U)$.
Also $R_{\rm nil}(U)\subset U$, so the intersection
of the residuals is trivial. A finite family of coset
conditions on the right can be combined in a finer
quotient. Compactness gives surjectivity and the
intersection calculation gives injectivity. Compatible
automorphisms and their compatible inverses act on
this inverse limit. This states the required transition
compatibility rather than assuming that arbitrary
members of the individual automorphism groups form
an inverse system. Requiring a continuous word
identity at every level is another closed condition;
its solvability is not implied by independent lifts.
\end{proof}

\subsection{The manuscript dictionary and dependency repairs}

\begin{statement}{B.8}{Proved}{The primitive, the cocycle, and the numerical coordinate are distinct}
In the normalized multiplier-one modular action,
\[
 \Phi(r)=ArA^{-1}=f^{-1}r
 \quad\Longrightarrow\quad
 f^{-1}=A\rho(A)^{-1},\qquad \rho=\operatorname{Ad}(r)|_F.
\]
Thus the manuscript's $g$ in
$u\rho=\operatorname{int}(g)\rho u$ is a cocycle,
and $A$ is its primitive. It is not the same object
as the Gaussian primitive $\beta$, nor as the scalar
$\rho_\Phi$ defined by
$(\beta^{-1})^{\rm ab}=(\rho_\Phi,\rho_\Phi)$.
\end{statement}
\begin{proof}
Multiply the displayed conjugation equation on the
right by $r^{-1}$. This gives
$f^{-1}=ArA^{-1}r^{-1}=A\rho(A)^{-1}$.
The three objects have different types: $g,A,\beta$
are group elements, whereas $\rho_\Phi$ is a profinite
integer. The cocycle and a primitive cannot be
interchanged in a finite evaluation. Cote 145,
scans 7, 11--12, 22 and 28, is available here only
through the supplied first-pass transcription;
the exact dictionary for the manuscript's separate
letter $\beta$ requires its written relation and
cannot be inferred merely from a modern reuse of
the same letter. We do not identify it with
$\rho_\Phi$ on typographical grounds.
\end{proof}

\begin{statement}{B.9}{Proved}{The normalized pair determines the outer gentle element}
In the fixed free generators and Belyi normalization,
the pair $(\lambda,f)$ with $f\in F'$ determines
an outer gentle element injectively. This statement
does not say that its residue modulo 12 determines
the two primitive frames.
\end{statement}
\begin{proof}
The pair specifies $x\mapsto x^\lambda$ and
$y\mapsto f^{-1}y^\lambda f$, hence the automorphism.
Conversely two such normalized representatives of
the same outer action differ by conjugation with
an element centralizing $x^\lambda$, hence by a
power $x^a$. Comparing the $y$ images shows that
the two conjugators differ by a power of $y$ and
this power of $x$. Their zero abelianizations
force both powers to vanish; centralizers of
nontrivial generator powers in a free profinite
group are the corresponding procyclic subgroups.
Thus the normalized $f$ is unique. A.9 gives a
specific reason why a small finite action can lose
the primitive-frame information.
\end{proof}

\begin{statement}{B.10}{Proved}{The corrected dependency order used in this return}
Every reference here to the CM field statement is
written ``LM7 D.25'', rather than the ambiguous
``D.25''. The R4 input is the corrected E.65 proof
including the pure-frame existence argument and
the explicit two-model Galois sign comparison.
The lifting and weight proofs in Part P and B.20--B.22
are ordered without the circular dependencies
E.2/E.58, D.1/D.51 and D.2/D.52.
\end{statement}
\begin{proof}
This is a statement about the dependencies of this
text, which can be checked from its displayed
proofs. The primitive existence uses the torsion
conjugacy and centralizer theorem, specifically
Herfort--Ribes (1985), Theorem B$'$, and the
congruence convention uses LM7 D.15. The explicit
field calculation appears independently in A.13;
E.62 is not an input. The positive residue in R2
is obtained from the relative lifting theorem
before it is used to prove that the sign is nonzero.
The weight computations are integrated before
their consequences are stated. The Hopficity
argument is the finite-open-subgroup permutation
argument in P.10, not an implication from finite
surjectivity without injectivity. No use of Turaev
2009 is made. These repairs concern the present
paper and its register; they do not silently
rewrite the supplied archive.
\end{proof}
